\section{Low-$F$ Multistandarding}
\label{app:low-f}

The Main theorem deliberately compares the high- and intermediate-cost regions. For completeness, this appendix records the verified sufficient low-$F$ region in the symmetric model. The result is not the selective-erosion mechanism: once every firm supports every relevant standard, the formal regimes differ only through duplicated fixed adoption costs.

Suppose
\begin{equation}
0<F<\min\{T_U,T_A\}.
\label{eq:appendix-low-f-condition}
\end{equation}
Because $T_W>T_U$, condition \eqref{eq:appendix-low-f-condition} also implies $F<T_W$. Under an SU, the outsider strictly adopts the bloc standard because $F<T_A<2T_A$. For the two members, adoption of the outsider standard is strictly profitable both when the other member has not adopted, since $F<T_U$, and when the other member has adopted, since $F<T_A$. Thus all three firms support both formal standards.

Under SW, for every target standard the two foreign firms strictly prefer adoption whether the rival has adopted or not: $F<T_W$ and $F<T_A$. Each firm therefore supports both foreign standards in addition to its native standard. International standardization already has a single common standard and requires no private adoption.

Universal multistandarding makes every destination a complete-compatibility product market. Consumer surplus and gross operating profits are therefore the same as under IS. The only remaining difference is the number of additional-standard fixed costs paid by each domestic firm. Under an SU each firm supports one additional standard, whereas under SW each firm supports two. Hence, for every country,
\begin{equation}
\boxed{
W_i^{SU}=W_i^{IS}-F,
\qquad
W_i^{SW}=W_i^{IS}-2F.
}
\label{eq:appendix-low-f-welfare}
\end{equation}
Since $F>0$, every country receives a lower continuation payoff under either non-IS formal regime than under IS in this sufficient low-cost region. This is the frozen low-$F$ formal-regime comparison recorded in the theorem ledger. We do not use it to extend Theorem~\ref{thm:formal-coalition-stability} beyond its stated high- and intermediate-$F$ regions.

This Appendix result should not be extrapolated into a global monotonicity statement in $F$. Its economics differs from the headline intermediate-cost result. In the intermediate region, outsider-only adoption selectively removes the regional exclusion rent while reciprocal incompatibility survives. In the sufficient low-$F$ region, effective compatibility is already universal and the remaining formal-regime difference is duplicate private adoption cost accounting. The Main contribution therefore remains the selective-erosion transition established in Sections~\ref{sec:selective-erosion} and \ref{sec:coalition-stability}.
